\documentclass[12pt, a4paper]{article}

\usepackage[utf8]{inputenc}
\usepackage[T1]{fontenc}
\usepackage{amsmath, amssymb, amsthm}
\usepackage{mathtools}
\usepackage{booktabs}
\usepackage{graphicx}
\usepackage{hyperref}
\usepackage[margin=2.5cm]{geometry}
\usepackage{natbib}
\usepackage{caption}
\usepackage{subcaption}
\usepackage{enumitem}

\newtheorem{theorem}{Theorem}
\newtheorem{proposition}[theorem]{Proposition}
\newtheorem{definition}{Definition}
\newtheorem{remark}{Remark}

\title{A Replicator--Mutator Framework for Ideological Dynamics:\\
Dormancy, Martyrdom, and Endogenous Landscape Modification}

\author{%
  Working Paper --- March 2026
}

\date{}

\begin{document}

\maketitle

\begin{abstract}
We present a mathematical framework for modeling the dynamics of competing
ideological systems in human populations. Standard evolutionary game theory
treats competing strategies on a fixed fitness landscape with absorbing
extinction boundaries --- assumptions that fail for ideas, which persist in
latency, gain strength from suppression, hybridize to form new species, and
modify the very environment in which they compete. Our framework extends the
replicator equation in three directions: (1)~a latent sympathy variable that
models dormancy and martyrdom dynamics, yielding a replicator--mutator system
where extinction is no longer absorbing; (2)~endogenous environment
modification, where dominant ideologies reshape the fitness landscape for all
competitors; and (3)~hybridization as a stochastic jump process, enabling the
emergence of new ideological species. The resulting system is a
piecewise-deterministic Markov process (PDMP). We implement this framework in
a simulation engine and apply it to a scenario of four competing ideological
currents in Europe from 1800 to 2025, with exogenous historical shocks
calibrated to major crises, wars, and technological disruptions. The
simulation reproduces qualitative features of ideological history including
the persistence of traditionalism despite modernization, the crisis-driven
surges of radical movements, and the dormancy--reawakening cycle. We discuss
the model's structural limitation: it treats individuals as passive hosts
rather than deliberate agents, a constraint we regard as both the source of
its tractability and the boundary of its explanatory reach.
\end{abstract}

\section{Introduction}
\label{sec:introduction}

Ideas compete. This observation is as old as intellectual history itself, but
its formal treatment has remained surprisingly underdeveloped. While
evolutionary biology has produced a rich mathematical apparatus for modeling
competition between species and strategies
\citep{hofbauer1998evolutionary,nowak2006evolutionary,weibull1995evolutionary},
and while cultural evolution theory has adapted some of these tools to the
study of beliefs and norms
\citep{cavalli1981cultural,boyd2005origin,mesoudi2011cultural}, the
specific dynamics of \emph{ideological} competition --- the struggle between
large-scale systems of political, religious, and philosophical thought ---
has largely been left to historians and political theorists working without
formal models.

This is not for lack of metaphor. Since \citet{dawkins1976selfish} introduced
the concept of the meme, the analogy between biological and cultural
evolution has been widely explored. But the analogy breaks in places that
matter. Three structural differences between species and ideas demand
mathematical attention:

\begin{enumerate}[label=(\roman*)]
  \item \textbf{Dormancy.} Ideas do not die when their carriers die or
  abandon them. Christianity survived the fall of Rome. Marxism survived the
  collapse of the Soviet Union. An idea can persist in latency --- in
  libraries, in diaspora communities, in private sympathy --- and reawaken
  when conditions change. In the standard replicator equation, the boundary
  $x_i = 0$ is absorbing: a strategy driven to zero stays at zero forever.
  For ideas, this is empirically wrong.

  \item \textbf{Hybridization.} A rabbit cannot also be a fox, but a person
  can be a Christian Marxist, or a nationalist who believes in universal
  rights. When ideas cohabit in a population, they can merge to form new
  species: social democracy emerged from the synthesis of Marxist analysis
  and liberal institutions \citep{berman2006primacy}; liberation theology
  from Christianity and anti-colonial politics. Standard population dynamics
  does not accommodate speciation during competition.

  \item \textbf{Endogenous landscape modification.} The fitness of a
  biological species depends on a niche that changes slowly relative to
  evolutionary timescales. The fitness of an idea depends on a social
  environment --- institutions, economic conditions, media structures ---
  that the idea itself modifies when it achieves dominance. Liberalism
  creates markets; markets create the inequality that fuels anti-liberal
  reaction. The fitness landscape is not exogenous: it is partly constructed
  by the competitors themselves.
\end{enumerate}

These are not minor perturbations of ecological dynamics. They are
qualitatively different phenomena that produce qualitatively different
behavior: oscillations driven by self-undermining dominance, resurrection
from apparent extinction, and the creation of genuinely new competitors
during the struggle.

The present paper develops a formal framework that addresses all three
features. We extend the replicator equation
\citep{taylor1978evolutionary} in three corresponding directions.
First, we introduce a \emph{latent sympathy} variable $y_i(t)$ alongside the
manifest adoption share $x_i(t)$, producing a replicator--mutator system
\citep{komarova2004replicator,page2002unifying} in which latent sympathy
prevents absorbing extinction and enables re-emergence. Suppression of an
ideology increases its latent sympathy (the \emph{martyrdom effect}),
formalizing the historical observation that persecuted ideas gain moral
authority. Second, we couple the replicator dynamics to an environment vector
$e(t)$ whose evolution depends on the current ideological distribution,
producing a co-evolutionary system with non-stationary fitness landscapes.
Third, we overlay a stochastic hybridization process that creates new
ideological types when latent sympathy overlap exceeds a threshold.

The resulting system is a piecewise-deterministic Markov process (PDMP)
\citep{davis1995piecewise}: continuous replicator--mutator dynamics
punctuated by discrete stochastic events (hybridization, mutation, charismatic
transmission). We implement this system in a simulation engine and apply it
to a historical scenario: four ideological currents --- liberalism, Marxism,
religious traditionalism, and nationalism --- competing in Europe from 1800 to
2025, subject to exogenous shocks calibrated to major historical crises.

The simulation is not a prediction engine. It is a tool for testing whether a
set of hypothesized mechanisms, when combined and allowed to interact, can
reproduce the qualitative shape of ideological history. We find that the
framework reproduces several robust features: the persistence of
traditionalism despite sustained modernization pressure, the crisis-driven
surges of radical movements \citep{funke2016going}, the dormancy and
reawakening of suppressed ideas, and the slow ideological drift produced by
endogenous environment modification.

We are explicit about the framework's deepest limitation. The replicator
equation, even in its extended form, treats individuals as passive carriers
of ideas whose adoption probabilities are determined by fitness and
environment. It has no room for genuine deliberation --- for a person who
sits quietly and changes their mind because they have thought carefully.
This is not an oversight; it is a structural choice that makes the population-level
dynamics tractable. But it means the model describes the ecology of ideas
\emph{as seen from above}. The view from inside --- from the deliberating
individual --- requires a different mathematics, perhaps one grounded in
decision theory or phenomenology rather than population dynamics. We regard
this boundary as the most important open problem in the field.

The paper is organized as follows. Section~\ref{sec:related} surveys related
work in cultural evolution, memetics, and computational social science.
Section~\ref{sec:framework} presents the mathematical framework.
Section~\ref{sec:software} describes the simulation software.
Section~\ref{sec:experiments} details the experimental design --- the choice
of ideological currents, parameter calibration, and historical shock
schedule. Section~\ref{sec:results} analyzes the simulation results.
Section~\ref{sec:discussion} discusses the contributions, limitations, and
implications of the framework.

\section{Related Work}
\label{sec:related}

Our framework draws on and departs from three research traditions: cultural
evolution theory, memetics, and computational models of opinion dynamics.

\subsection{Cultural Evolution}

The mathematical study of cultural transmission begins with
\citet{cavalli1981cultural}, who adapted population genetics models to
cultural traits, distinguishing vertical (parent--child), horizontal
(peer--peer), and oblique (elder--youth) transmission modes.
\citet{boyd2005origin} extended this program with models of biased
transmission, conformist bias, and prestige bias, showing that cultural
evolution can produce group-level adaptations not reducible to individual
selection. \citet{mesoudi2011cultural} provides a comprehensive synthesis.

These models treat cultural traits as discrete variants propagating through
a population --- formally, alleles in a cultural genome. The framework is
powerful for small-scale traits (tool designs, food taboos, phonemes) but
encounters difficulties with \emph{ideologies}: large, internally structured
systems of belief that interact with institutions, respond to economic
conditions, and actively reshape the environments in which they propagate.
The cultural evolution tradition typically assumes a fixed or slowly varying
selective environment; our framework makes the environment endogenous.

\citet{bisin2001economics} developed an economic model of cultural
transmission where parents invest in socializing children into their own
cultural group. Their model produces long-run coexistence of cultural types
through a mechanism analogous to our latent sympathy: minority groups invest
more heavily in transmission, preventing assimilation. Our latent variable
$y_i$ generalizes this insight beyond the parent--child dyad.

\subsection{Memetics}

\citet{dawkins1976selfish} proposed that ideas (``memes'') replicate, mutate,
and compete in a manner analogous to genes. The analogy has been productive
as a heuristic but has resisted formalization. \citet{shennan2002genes}
applied memetic reasoning to archaeological data with some success, but the
field has struggled with a core problem: unlike genes, ideas do not have a
well-defined unit of replication. A political ideology is not a single meme
but a complex of mutually reinforcing propositions, aesthetic preferences,
institutional practices, and identity markers.

Our framework sidesteps this problem by working at the level of entire
ideological systems rather than individual memes. We model ideologies as
indivisible competitors on a simplex --- each characterized by a vector of
responses to environmental conditions --- rather than attempting to decompose
them into memetic atoms. This is a deliberate coarse-graining: it sacrifices
internal structure for tractability at the population level.

\subsection{Opinion Dynamics and Social Contagion}

A large literature in computational social science models opinion spread as a
contagion or threshold process. \citet{granovetter1978threshold} introduced
threshold models where individuals adopt a behavior when a sufficient
fraction of their neighbors have adopted it. \citet{watts1998collective}
showed that small-world network structure can produce cascading adoption.
\citet{acemoglu2011bayesian} studied Bayesian learning in social networks,
showing how network structure determines whether populations converge to
correct beliefs.

These models typically operate at the individual level and model opinion as a
continuous variable or binary choice. Our framework operates at the
population level and models competition between multiple discrete ideologies.
The two approaches are complementary: network models explain the
micro-mechanisms of spread (why \emph{this} person adopted \emph{this} idea
now), while our framework explains the macro-dynamics (why liberalism grew
in the 19th century and contracted in the 21st).

\subsection{Empirical Studies of Ideological Change}

The empirical literature provides the calibration targets for our model.
\citet{funke2016going} analyzed 800 elections in 20 democracies over 140
years and found that financial crises increase far-right vote share by
approximately 30\%. \citet{inglehart1997modernization} documented the
``post-materialist'' value shift across generations using 50 years of World
Values Survey data, confirming the generational transmission mechanism.
\citet{jennings1968transmission} showed that parental party identification
predicts children's at $r = 0.47$--$0.59$. \citet{pierson2000increasing}
formalized the ``policy ratchet'' --- the tendency of institutional changes
to become self-reinforcing --- providing the empirical basis for our
institutional amplification mechanism. \citet{norris2019cultural}
documented the rise of authoritarian populism as a cultural backlash against
post-materialist values, a dynamic our model captures through the
interaction between environmental variables and ideological fitness.

\citet{king2008ordinary} showed that unemployment correlated with radical
vote share across Weimar Germany's electoral districts, providing perhaps the
cleanest quantitative evidence for crisis recruitment. These findings ground
our treatment of exogenous shocks as temporary increases in the
susceptibility of populations to ideological change.

\subsection{What Is Missing}

No existing framework, to our knowledge, combines all three features we
address: dormancy with martyrdom dynamics, endogenous environment
modification, and hybridization as speciation. Cultural evolution models
handle transmission but not institutional feedback. Memetic models handle
competition but not latency. Opinion dynamics models handle micro-level
spread but not macro-level ideological competition. Our contribution is to
integrate these mechanisms into a single, mathematically tractable system.

\input{sec_framework}
\input{sec_software}
\input{sec_experiments}
\input{sec_results}
\input{sec_discussion}

\bibliographystyle{apalike}
\bibliography{references}

\end{document}
